% ---------------------------------------------------------------------------
% BFA-4R: the evidence a publication supplied.
%
% Paste this fragment into the manuscript. It needs booktabs, array, tikz.
% It defines \gf, \gh and \go with \providecommand, so a definition already in
% the manuscript wins and this one is ignored.
%
% This axis is deliberately SEPARATE from the recomputation status of the
% catalog table. BFA-4R grades the form and quality of the computational
% evidence the ORIGINAL PUBLICATION made available. The recomputation status
% grades what THIS work verified independently. They are not interchangeable,
% and merging them into one status would destroy the comparison.
%
% Rows are generated by manuscript/generate_bfa4r_table.jl.
% ---------------------------------------------------------------------------

% Harvey balls. Greyscale-safe: the distinction is fill geometry, not colour.
% Robust: these appear inside \caption, whose contents are written to the .aux
% list-of-tables entry, and a fragile tikzpicture there breaks the build.
\providecommand{\gball}{}
\DeclareRobustCommand{\gball}[1]{%
  \begin{tikzpicture}[baseline=-0.55ex,scale=0.85]
    \draw[line width=0.4pt] (0,0) circle (0.52ex);
    \ifnum#1>0 \fill (0,0) -- (90:0.52ex) arc (90:{90-180*#1}:0.52ex) -- cycle; \fi
  \end{tikzpicture}}
\providecommand{\gf}{}\DeclareRobustCommand{\gf}{\gball{2}}  % full  — criterion met
\providecommand{\gh}{}\DeclareRobustCommand{\gh}{\gball{1}}  % half  — partly met
\providecommand{\go}{}\DeclareRobustCommand{\go}{\gball{0}}  % open  — not met

\subsection*{BFA-4R: the evidence the publication supplied}

BFA-4R grades a \emph{publication}, not a family, a construction, a catalog
entry or a claim. It asks what computational evidence the paper itself made
available, and it is independent of whether that evidence turned out to be
correct. A paper can supply exemplary evidence for a claim that later
recomputation contradicts, and a paper can supply nothing for a claim that
recomputation confirms. Both cases occur in this collection, which is why the
recomputation status of the catalog table and the BFA-4R grade below are
reported on separate axes and are never combined into a single score.

\begin{description}\itemsep2pt
\item[R1 -- Methodological Disclosure.] Captures whether a paper could be
reimplemented from its own text. \gf\ when the paper provides explicit
pseudocode or a complete parameter specification (exact field, primitive
polynomial, matrix structure, search bounds, runtime) sufficient for clean-room
reimplementation without guesswork. \gh\ when the mathematical structure and key
formulas are given, but search heuristics, pruning strategies, or equivalence
methods are underspecified. \go\ when the method is described only in
qualitative terms, with no explicit algorithms or parameters.

\item[R2 -- Numerical Verification Metadata.] Captures whether a paper's central
claim can be checked against its own reported outputs without rerunning code.
\gf\ when the paper provides a systematic table of invariants ($N_F$-sequences,
$\Gamma$/$\Delta$-ranks, spectra) across multiple dimensions or function
families, sufficient to verify the main claim. \gh\ when only sparse invariant
evidence is present, such as a few key values cited in-text or a table covering
a single dimension. \go\ when the computational claim is asserted with no
supporting output.

\item[R3 -- Artifact Availability.] Captures whether code or data is shared in a
form that survives and can be reused. \gf\ when an open-source artifact
(executable code with resolved dependencies, or a versioned dataset on a
persistent archive such as Zenodo or a citable GitHub release) is publicly
available and sufficient to reuse the core output. \gh\ when a public artifact
is incomplete, proprietary-dependent (e.g., a Magma script), released without
its generation code, or hosted on a transient personal page. \go\ when no
computational artifact is publicly available.

\item[R4 -- External Corroboration.] Captures whether anyone outside the author
team has confirmed the result. \gf\ when a named third party explicitly confirms
the core computational claim using a different method, with proper citation.
\gh\ when independent events partially support the claim, such as
cross-referencing in a subsequent comparison table, rediscovery via a different
method, or inclusion in a community database. \go\ when all validation is purely
internal to the authors, with no external reference points.
\end{description}

\begin{table}[htbp]
\centering
\caption[BFA-4R grades of the reproduced publications]{BFA-4R grades for the \bfaPapers{} publications reproduced in this
work, ordered by year. One row is one \emph{publication}. \gf\ criterion met,
\gh\ partly met, \go\ not met; the marks differ by fill geometry and carry no
colour. R1 methodological disclosure, R2 numerical verification metadata, R3
artifact availability, R4 external corroboration. R3 is met in full by
\bfaFullRthree{} of \bfaPapers{} publications and R4 by \bfaFullRfour{}, against
\bfaFullRtwo{} for R2. \emph{These grades are rule-derived proposals, not an
expert reading.} Each is produced by a stated rule over the paper source and the
reproduction review, and the raw evidence behind every cell is given in
Table~\ref{tab:bfa4r-evidence} so the reader can audit and override it. The
paper sources are automated PDF-to-LaTeX conversions in which the LaTeX \texttt{algorithm} and \texttt{theorem}
environments occur zero times across all \bfaPapers{} papers, so R1 rests on a textual proxy and is the least reliable
column after R4. \textsuperscript{\dag}~no usable bibliography key
(\bfaNoBib{} publication). This table grades the evidence a publication
supplied; it does not grade whether that evidence is correct, which is the
separate recomputation axis.}
\label{tab:bfa4r}
\small
\begin{tabular}{@{}l l cccc @{\hspace{2em}} l l cccc@{}}
\toprule
Publication & Ref. & R1 & R2 & R3 & R4 &
Publication & Ref. & R1 & R2 & R3 & R4 \\
\midrule
% !!! GENERATED FILE — DO NOT EDIT !!!
% Produced by manuscript/generate_bfa4r_table.jl from:
%   papers/<year>/<id>/input/*.tex     (the publication, PDF-converted)
%   papers/<year>/<id>/README.md       (the reproduction review)
%   packages/VBFCatalog/data-raw/bibliography.bib
%   packages/VBFCatalog/BFA_catalog.sqlite  (bulk_corpora, functions)
% Regenerate with:  julia --project=. manuscript/generate_bfa4r_table.jl
\texttt{\scriptsize Biryukov\_200300} & \cite{Biryukov_200300} & \gf & \gf & \go & \gh & \texttt{\scriptsize Yu\_202012} & \cite{Yu_202012} & \gf & \gh & \gh & \gh \\
\texttt{\scriptsize Budaghyan\_200506} & \cite{Budaghyan_200506} & \go & \gf & \go & \gh & \texttt{\scriptsize Beierle\_202107} & \cite{Beierle_202107} & \gf & \gh & \gf & \go \\
\texttt{\scriptsize Bracken\_200807} & \cite{Bracken_200807} & \gh & \gh & \gh & \gh & \texttt{\scriptsize Budaghyan\_202101} & \cite{Budaghyan_202101} & \gh & \gf & \gh & \gh \\
\texttt{\scriptsize Brinkmann\_200812} & \cite{Brinkmann_200812} & \gf & \gf & \gh & \gh & \texttt{\scriptsize Calderini\_202100} & \cite{Calderini_202100} & \go & \gf & \gh & \gh \\
\texttt{\scriptsize Budaghyan\_200904} & \cite{Budaghyan_200904} & \go & \gh & \go & \gh & \texttt{\scriptsize Chase\_202110} & \cite{Chase_202110} & \gh & \gh & \gh & \gh \\
\texttt{\scriptsize Budaghyan\_200910} & \cite{Budaghyan_200910} & \gh & \gh & \go & \go & \texttt{\scriptsize Beierle\_202201} & \cite{Beierle_202201} & \gf & \gh & \gf & \gf \\
\texttt{\scriptsize Edel\_200900} & \cite{Edel_200900} & \gh & \gf & \gh & \gf & \texttt{\scriptsize Canteaut\_202205} & \cite{Canteaut_202205} & \gf & \gf & \gh & \gh \\
\texttt{\scriptsize Bracken\_201103} & \cite{Bracken_201103} & \gh & \gf & \gh & \gh & \texttt{\scriptsize Gologlu\_202207} & \cite{Gologlu_202207} & \gf & \gf & \gh & \gh \\
\texttt{\scriptsize Bracken\_201205} & \cite{Bracken_201205} & \gh & \gh & \gh & \go & \texttt{\scriptsize Kaleyski\_202200} & \cite{Kaleyski_202200} & \go & \gh & \go & \go \\
\texttt{\scriptsize Yuyin\_201300} & \cite{Yuyin_201300} & \gf & \gh & \go & \gh & \texttt{\scriptsize Kaleyski\_202203} & \cite{Kaleyski_202203} & \gh & \gf & \gh & \go \\
\texttt{\scriptsize Li\_201408} & \cite{Li_201408} & \go & \gh & \gh & \go & \texttt{\scriptsize Li\_202207} & \cite{Li_202207} & \go & \gf & \gf & \gh \\
\texttt{\scriptsize Zha\_201401} & \cite{Zha_201401} & \go & \gf & \gh & \go & \texttt{\scriptsize Yu\_202211d} & \cite{Yu_202211} & \gf & \gh & \gh & \gf \\
\texttt{\scriptsize Tang\_201510} & \cite{Tang_201510} & \gh & \gf & \gh & \gh & \texttt{\scriptsize Zheng\_202211} & \cite{Zheng_202211} & \go & \gf & \go & \go \\
\texttt{\scriptsize Canteaut\_201906} & \cite{Canteaut_201906} & \gf & \gh & \go & \go & \texttt{\scriptsize Calderini\_202306} & \cite{Calderini_202306} & \gh & \gf & \gh & \gh \\
\texttt{\scriptsize Kaleyski\_201911} & \cite{Kaleyski_201911} & \go & \gf & \go & \go & \texttt{\scriptsize Ichanska\_202400} & \cite{Ichanska_202400} & \gf & \gf & \gh & \gh \\
\texttt{\scriptsize Taniguchi\_201909} & \cite{Taniguchi_201909} & \gh & \gh & \gh & \gh & \texttt{\scriptsize Li\_202402} & \cite{Li_202402} & \gh & \gf & \go & \go \\
\texttt{\scriptsize Budaghyan\_202008} & \cite{Budaghyan_202008} & \gh & \gh & \gh & \go & \texttt{\scriptsize Beierle\_202508} & \cite{Beierle_202508} & \gh & \gf & \gf & \gh \\
\texttt{\scriptsize Budaghyan\_202009} & \cite{Budaghyan_202009} & \gf & \gf & \gh & \go & \texttt{\scriptsize Gologlu\_202503} & \textsuperscript{\dag} & \gh & \gf & \gh & \gh \\
\texttt{\scriptsize Budaghyan\_202011} & \cite{Budaghyan_202011} & \gh & \gf & \gh & \gh & \texttt{\scriptsize Taniguchi\_202501} & \cite{Taniguchi_202501} & \go & \gh & \gh & \gh \\
\texttt{\scriptsize Calderini\_202009} & \cite{Calderini_202009} & \gf & \gf & \gh & \gh & \texttt{\scriptsize Yu\_202512} & \cite{Yu_202512} & \gf & \gh & \gf & \gh \\
\texttt{\scriptsize Kalgin\_202000} & \cite{Kalgin_202000} & \gf & \gf & \gh & \go & \texttt{\scriptsize Kuznetsov\_202606} & \cite{Kuznetsov_202600} & \gf & \gf & \gf & \gh \\

\bottomrule
\end{tabular}
\end{table}

\clearpage

\begin{table}[htbp]
\centering
\caption[BFA-4R raw evidence]{Raw evidence behind every cell of Table~\ref{tab:bfa4r}, so each grade
can be audited. \textsc{pc} textual pseudocode markers; \textsc{io}
\texttt{Input:}/\texttt{Output:}/\texttt{Require:}/\texttt{Ensure:} lines;
\textsc{par} explicit-parameter and cost mentions; \textsc{tab}
\texttt{tabular} environments; \textsc{inv} invariant mentions; \textsc{arc}
persistent-archive references; \textsc{rep} code-repository references;
\textsc{pro} proprietary-tool references; \textsc{und} underspecification
findings recorded in the reproduction review's \emph{Paper Issues} section;
\textsc{cor} other reproductions in this collection naming this paper;
\textsc{cps} the paper's functions appear in a bulk corpus. R1 uses
\textsc{pc}, \textsc{io}, \textsc{par} and \textsc{und}; R2 uses \textsc{tab}
and \textsc{inv}; R3 uses \textsc{arc}, \textsc{rep} and \textsc{pro}; R4 uses
\textsc{cor} and \textsc{cps}.}
\label{tab:bfa4r-evidence}
\scriptsize
\begin{tabular}{@{}l *{11}{r}@{}}
\toprule
Publication & \textsc{pc} & \textsc{io} & \textsc{par} & \textsc{tab} &
\textsc{inv} & \textsc{arc} & \textsc{rep} & \textsc{pro} & \textsc{und} &
\textsc{cor} & \textsc{cps} \\
\midrule
% !!! GENERATED FILE — DO NOT EDIT !!!
% Produced by manuscript/generate_bfa4r_table.jl from:
%   papers/<year>/<id>/input/*.tex     (the publication, PDF-converted)
%   papers/<year>/<id>/README.md       (the reproduction review)
%   packages/VBFCatalog/data-raw/bibliography.bib
%   packages/VBFCatalog/BFA_catalog.sqlite  (bulk_corpora, functions)
% Regenerate with:  julia --project=. manuscript/generate_bfa4r_table.jl

\texttt{Biryukov\_200300} & 80 & 0 & 11 & 5 & 11 & 0 & 0 & 0 & 0 & 1 & n \\
\texttt{Budaghyan\_200506} & 0 & 0 & 0 & 3 & 25 & 0 & 0 & 0 & 0 & 3 & n \\
\texttt{Bracken\_200807} & 8 & 0 & 0 & 0 & 13 & 0 & 0 & 3 & 0 & 2 & n \\
\texttt{Brinkmann\_200812} & 17 & 0 & 5 & 7 & 23 & 0 & 0 & 5 & 0 & 3 & n \\
\texttt{Budaghyan\_200904} & 0 & 0 & 0 & 0 & 22 & 0 & 0 & 0 & 0 & 1 & n \\
\texttt{Budaghyan\_200910} & 2 & 0 & 0 & 0 & 8 & 0 & 0 & 0 & 0 & 0 & n \\
\texttt{Edel\_200900} & 0 & 0 & 1 & 12 & 46 & 0 & 0 & 6 & 0 & 3 & y \\
\texttt{Bracken\_201103} & 5 & 0 & 0 & 3 & 25 & 0 & 0 & 8 & 0 & 1 & n \\
\texttt{Bracken\_201205} & 1 & 0 & 0 & 1 & 35 & 0 & 0 & 1 & 0 & 0 & n \\
\texttt{Yuyin\_201300} & 1 & 0 & 1 & 1 & 31 & 0 & 0 & 0 & 0 & 2 & n \\
\texttt{Li\_201408} & 0 & 0 & 0 & 0 & 63 & 0 & 0 & 3 & 0 & 0 & n \\
\texttt{Zha\_201401} & 0 & 0 & 0 & 2 & 48 & 0 & 0 & 4 & 0 & 0 & n \\
\texttt{Tang\_201510} & 1 & 0 & 0 & 4 & 93 & 0 & 0 & 11 & 0 & 1 & n \\
\texttt{Canteaut\_201906} & 1 & 0 & 1 & 0 & 28 & 0 & 0 & 0 & 0 & 0 & n \\
\texttt{Kaleyski\_201911} & 0 & 0 & 0 & 2 & 15 & 0 & 0 & 0 & 0 & 0 & n \\
\texttt{Taniguchi\_201909} & 0 & 0 & 2 & 1 & 4 & 0 & 0 & 3 & 0 & 0 & y \\
\texttt{Budaghyan\_202008} & 0 & 0 & 1 & 1 & 19 & 0 & 0 & 3 & 0 & 0 & n \\
\texttt{Budaghyan\_202009} & 6 & 0 & 4 & 3 & 34 & 0 & 0 & 4 & 0 & 0 & n \\
\texttt{Budaghyan\_202011} & 1 & 0 & 0 & 5 & 48 & 0 & 0 & 1 & 0 & 1 & n \\
\texttt{Calderini\_202009} & 4 & 0 & 1 & 5 & 22 & 0 & 0 & 3 & 0 & 1 & n \\
\texttt{Kalgin\_202000} & 6 & 0 & 1 & 9 & 43 & 0 & 0 & 3 & 0 & 0 & n \\
\texttt{Yu\_202012} & 6 & 2 & 4 & 1 & 27 & 0 & 0 & 1 & 0 & 0 & y \\
\texttt{Beierle\_202107} & 34 & 2 & 12 & 1 & 29 & 2 & 2 & 2 & 0 & 0 & n \\
\texttt{Budaghyan\_202101} & 1 & 0 & 0 & 4 & 24 & 0 & 0 & 3 & 0 & 1 & n \\
\texttt{Calderini\_202100} & 0 & 0 & 0 & 6 & 65 & 0 & 0 & 2 & 1 & 5 & n \\
\texttt{Chase\_202110} & 3 & 0 & 0 & 0 & 8 & 0 & 2 & 0 & 0 & 1 & n \\
\texttt{Beierle\_202201} & 23 & 2 & 7 & 1 & 60 & 3 & 2 & 2 & 0 & 1 & y \\
\texttt{Canteaut\_202205} & 76 & 0 & 24 & 8 & 79 & 0 & 6 & 9 & 0 & 2 & n \\
\texttt{Gologlu\_202207} & 3 & 0 & 2 & 3 & 30 & 0 & 0 & 1 & 0 & 2 & n \\
\texttt{Kaleyski\_202200} & 0 & 0 & 0 & 0 & 20 & 0 & 0 & 0 & 0 & 0 & n \\
\texttt{Kaleyski\_202203} & 20 & 6 & 15 & 3 & 35 & 0 & 0 & 5 & 1 & 0 & n \\
\texttt{Li\_202207} & 0 & 0 & 0 & 5 & 23 & 2 & 4 & 11 & 0 & 1 & n \\
\texttt{Yu\_202211d} & 11 & 0 & 7 & 0 & 28 & 0 & 2 & 4 & 0 & 4 & y \\
\texttt{Zheng\_202211} & 0 & 0 & 0 & 4 & 54 & 0 & 0 & 0 & 1 & 0 & n \\
\texttt{Calderini\_202306} & 1 & 0 & 0 & 8 & 24 & 0 & 0 & 1 & 1 & 4 & n \\
\texttt{Ichanska\_202400} & 34 & 0 & 15 & 13 & 32 & 0 & 2 & 4 & 0 & 2 & n \\
\texttt{Li\_202402} & 2 & 0 & 0 & 7 & 110 & 0 & 0 & 0 & 0 & 0 & n \\
\texttt{Beierle\_202508} & 8 & 1 & 1 & 2 & 23 & 5 & 2 & 0 & 3 & 2 & n \\
\texttt{Gologlu\_202503} & 2 & 0 & 0 & 3 & 35 & 0 & 0 & 2 & 0 & 1 & n \\
\texttt{Taniguchi\_202501} & 0 & 0 & 0 & 1 & 18 & 0 & 0 & 2 & 0 & 1 & n \\
\texttt{Yu\_202512} & 64 & 2 & 15 & 1 & 45 & 5 & 2 & 7 & 0 & 2 & n \\
\texttt{Kuznetsov\_202606} & 15 & 2 & 17 & 7 & 60 & 12 & 8 & 16 & 0 & 1 & n \\

\bottomrule
\end{tabular}
\end{table}
